\section{把现实选择翻译成优化问题}
\label{sec:09-Convex-Optimization/01-Optimization-Foundations/01}

需求文档上只有一句话：``把首页推荐做得更好一点。''

这句话没有一处是错的，也没有一处能直接交给求解器。它没说能动什么——是排序阶段的打分，还是每个位次放什么，还是连召回的内容池一起换掉；没说``更好''按什么记分——点击、停留时长、次日留存，还是三个月后的付费；更没说哪些事情不许发生——广告密度、单个作者的曝光上限、每次请求的推理预算。

要把它变成一个可求解的对象，得交出三样东西：\textbf{决策变量}，也就是我们能动的量；\textbf{目标函数}，一个用来给方案排序的数；\textbf{约束}，任何方案都不许违反的条件。这一部分后面所有的技术——凸集、对偶、内点法——都站在这三样已经写对的前提上。而写它们的时候没有算法可用，只有判断。

\begin{center}
\begin{tikzpicture}[x=1cm,y=1cm,font=\small,>=stealth]
  \node[font=\footnotesize,black!70] at (6.45,6.0) {写下来的部分};
  \node[font=\footnotesize,black!70] at (11.3,6.0) {同时塞进去的假设};
  \draw[fill=black!6,draw=black!50,rounded corners=2pt] (0,2.35) rectangle (3.2,3.85);
  \node[font=\footnotesize,align=center] at (1.6,3.1) {需求文档上的一句话\\``把首页推荐\\做得更好一点''};
  \draw[fill=blue!10,draw=black!50] (4.5,4.35) rectangle (8.4,5.55);
  \node[font=\footnotesize,align=center] at (6.45,4.95) {决策变量\\第 \(k\) 个位次展示哪条内容};
  \draw[fill=blue!10,draw=black!50] (4.5,2.5) rectangle (8.4,3.7);
  \node[font=\footnotesize,align=center] at (6.45,3.1) {目标函数\\最大化本次会话的点击总数};
  \draw[fill=blue!10,draw=black!50] (4.5,0.65) rectangle (8.4,1.85);
  \node[font=\footnotesize,align=center] at (6.45,1.25) {约束\\每屏广告 \(\le1\) 条，\\单作者曝光占比 \(\le20\%\)};
  \draw[->,black!60] (3.25,3.1) -- (4.45,4.95);
  \draw[->,black!60] (3.25,3.1) -- (4.45,3.1);
  \draw[->,black!60] (3.25,3.1) -- (4.45,1.25);
  \draw[dashed,fill=black!3,draw=black!45] (9.0,4.35) rectangle (13.6,5.55);
  \node[font=\footnotesize,align=center,black!75] at (11.3,4.95) {内容池、召回与生产端\\都当作固定不动};
  \draw[dashed,fill=black!3,draw=black!45] (9.0,2.5) rectangle (13.6,3.7);
  \node[font=\footnotesize,align=center,black!75] at (11.3,3.1) {点击等于满意，\\当期点击代表长期留存};
  \draw[dashed,fill=black!3,draw=black!45] (9.0,0.65) rectangle (13.6,1.85);
  \node[font=\footnotesize,align=center,black!75] at (11.3,1.25) {没写进来的都不重要：\\审核成本、多样性、推理预算};
  \draw[->,black!45] (8.45,4.95) -- (8.95,4.95);
  \draw[->,black!45] (8.45,3.1) -- (8.95,3.1);
  \draw[->,black!45] (8.45,1.25) -- (8.95,1.25);
\end{tikzpicture}

\emph{一句自然语言变成三元组的过程。左边是需求，中间是交给求解器的东西，右边是没人签字、却随之生效的假设。}
\end{center}

图里最右边一列才是这一步的难处。每一次把模糊说法收紧成符号，都顺手固定了一批本来可以商量的东西，而这些东西不会出现在模型里，也不会出现在求解器的日志里。变量选定了，没被写成变量的杠杆就此消失；目标选定了，一个具体的偏序被安到了所有方案头上；约束选定了，可行域的边界被画在了一个具体的位置上。求解器不会质疑其中任何一条，它只会在你划出的范围里，尽最大努力把你写下的那个数推到最大。

\subsection{目标写歪了，解会精确地钻那个空子}

一个内容平台把排序目标定成会话内点击总数。上线两周，点击率涨了近一成，人均阅读时长掉了，第三周开始次日留存松动。事后复盘并不费力：在这个目标下，夸张标题配一张诱导图是严格占优的策略，系统只是找到了它。目标函数的公式没写错，写错的是``更好''这个词的定义，而优化过程忠实地照着执行。

另一个例子更短。客服工单系统把座席的考核指标改成平均处理时长，两个月后这个数字下降了两成，同时重复来电率上升了三成——提前结单和转派工单在新目标下都是划算的。

这类事故有个共同的形状。优化的结果几乎总落在可行域的边界上，而``边界''包含所有你没想到、也没禁止的做法。写下目标函数，等于宣布一个方案排序；凡是在这个排序里更靠前的方案，迟早会被找到，无论它在别的维度上有多难看。管理学里管这叫指标被劫持，强化学习里管这叫奖励塑形失败，二者说的是同一件事：代理指标与真实意图之间的缝隙，会被优化过程系统地撑大，而不是被磨平。目标写歪了，解得再准也只是更精确地答错。

补救的方向有两个，效果不同。一个是往目标里加惩罚项，比如减去举报率、加上阅读完成率；另一个是把它写成硬约束，比如要求举报率不超过某个值。两者不等价：带有限权重的惩罚项允许用足够大的收益换来轻微违规，硬约束不允许。选哪一个，取决于那条线是可赎买的还是不可赎买的——这是业务判断，不是数学判断，但它会一路影响到后面用什么算法。

\subsection{漏掉的约束会在上线前一天自己出现}

约束这一栏的失误方式相反，问题出在写少。一个排序模型在离线指标上明显更好，压测时才发现单次请求的推理量涨到三倍，超了延迟预算；一个信贷模型效果不错，合规审查时才发现某个受保护属性上的通过率差异越了线。这两件事在数学上是同一回事——可行域画大了，最优点跑到了一块本不该到达的区域里，而模型内部没有任何东西会提示这件事。求解器只回答被写下来的问题。

反过来，约束写多、写死也有代价，那就是不可行：没有任何点能同时满足全部条件。现实中的不可行大多不是算法的失败，而是几个部门的要求互相冲突被第一次摆到了一起。这时该改的是需求，不是求解器的参数。能在建模阶段发现冲突，比上线后发现便宜得多。

\subsection{三样东西写全了，问题才立得住}

换一个小到可以手算的例子，看看写全之后是什么样。一家工厂生产两种产品，A 每件利润 3，B 每件利润 5；两道工序每天分别只有 100 和 90 个工时，做一件 A 各耗 \((2,1)\) 个工时，做一件 B 各耗 \((1,3)\) 个工时。令 \(x_A,x_B\) 为两种产品的日产量，模型是

\[
\begin{aligned}
\text{maximize}\quad &3x_A+5x_B\\
\text{subject to}\quad
&2x_A+x_B\le100,\\
&x_A+3x_B\le90,\\
&x_A,x_B\ge0.
\end{aligned}
\]

满足全部约束的点组成可行域 \(C\)，这里是四个半平面相交得到的多边形。目标函数的等值线 \(3x_A+5x_B=c\) 是一族平行直线，让 \(c\) 增大就是把这族线往右上推，最后一次还碰到多边形的位置就是最优解。两条资源约束的边界交于 \((x_A,x_B)=(42,16)\)，利润 \(206\)；其余顶点 \((0,0)\)、\((50,0)\)、\((0,30)\) 的利润分别是 \(0\)、\(150\)、\(150\)。线性目标在多边形上的最优值总能在某个顶点取到，所以 \((42,16)\) 就是答案。

这里还有一个容易被漏掉的栏目：变量的定义域。若产品不可分割，模型就必须显式写上 \(x_A,x_B\in\Z\)。本例的连续最优点碰巧是整数，一般情况下不会这么巧，而补上这一行之后问题的性质会变——线性规划有多项式时间算法，整数规划在最坏情形下是 NP-hard 的，具体分界见\hyperref[sec:09-Convex-Optimization/09-Complexity-and-Approximation/01]{``P、NP 与 NP-hard：给问题按计算难度分类''}。不能因为求解器接受小数，就默认半台机器是可行方案。

\subsection{标准形式只是一门共同语言}

为了让理论和软件有统一的接口，优化问题通常写成

\[
\begin{aligned}
\text{minimize}\quad & f_0(x)\\
\text{subject to}\quad
& f_i(x)\le0,\quad i=1,\ldots,m,\\
& h_j(x)=0,\quad j=1,\ldots,p.
\end{aligned}
\]

最大化效用 \(u(x)\) 改写成最小化 \(-u(x)\)，\(\ge\) 型约束两边乘 \(-1\)，都不改变可行域，也不改变方案之间的排序，所以是安全的变形。但另一些改写就没这么无辜：对目标取一个非单调的变换会重排方案；把等式两边平方会引入符号相反的额外解；删掉一条``看起来不会起作用''的约束，只在你猜对了极端情形时才安全；用样本平均替代最坏情形，改掉的是整个风险口径。标准化让公式变短，可它不负责保存问题的含义，那件事得由人来核对。

同样属于建模的还有量纲和缩放。如果目标里一项以元计、另一项以毫秒计，两项相加在数学上无碍，数值上却会让求解器寸步难行；同理，一个变量取值在 \(10^{-6}\) 附近、另一个在 \(10^{6}\) 附近时，条件数会先于算法崩掉。这些是在调用求解器之前就该处理干净的事。

\subsection{最优值和最优解可以不是一回事}

设可行域为 \(C\)，最优值定义为

\[
p^\star=\inf_{x\in C}f_0(x).
\]

用下确界而不是最小值，是因为那个最好的数值未必真有可行点能取到。若存在 \(x^\star\in C\) 使 \(f_0(x^\star)=p^\star\)，它才叫最优解。在开区间 \(x>0\) 上最小化 \(x\) 就是个现成的例子：目标可以无限接近零，\(p^\star=0\)，但 \(x=0\) 不可行，最优解不存在；换成闭区间 \(x\ge0\)，同一个目标在 \(x^\star=0\) 处取到。存在性的一个够用的保障来自极值定理：\(C\) 非空且紧、\(f_0\) 在 \(C\) 上连续，最优值必定在某点取得。紧性不是必要条件，但它把闭性、有界性、连续性三件事为什么重要说清楚了。

于是一个优化问题至少有四种归宿，而``解出来了''只是其中一种。它可能\textbf{不可行}，可行域为空；可能\textbf{下方无界}，目标能一路降下去，比如在整条实线上最小化 \(-x\)，这有时来自漏写的容量约束，有时确实反映模型允许无限套利；可能\textbf{最优值有限却取不到}，答案有下界，只是可行域不含那个极限点；也可能\textbf{存在最优解}，这时才轮得到讨论怎么找、唯不唯一。求解器还会报告数值精度不足、迭代超限、``无法区分不可行与无界''，这些状态各有各的含义，任何一个都不等于``已经求得最优''。

到这里，问题算是写下来了。接下来求解器会返回一个点，并告诉你它在那里停住了——它周围再没有更好的方向。这句话离``这是所有可行方案里最好的''还有多远，是下一节要处理的事。
